\section{Books, Voices, and the Accidents of Survival}
\label{sec:origen-books}

\subsection{A corpus made by more hands than one}

Ancient lists portray extraordinary output, but no secure total can be made
from them. Eusebius says that he catalogued Origen's works in his lost life of
Pamphilus and names many in \emph{Ecclesiastical History} 6.23--36. Jerome's
\emph{On Illustrious Men} 54 declines to repeat the list and points instead to
his Letter 33 to Paula. Titles overlap, works had multiple books, stenographic
sermons could be edited, and the lists are incomplete. Familiar
claims that Origen wrote two thousand, six thousand, or some other exact number
usually conflate units and later reports.

The material system is clearer than the total. Ambrose's money supported
shorthand writers who took turns, copyists who produced ordinary manuscripts,
and women skilled in calligraphy. Origen dictated, revised, preached, answered
questions, and sometimes wrote letters in his own name; collaborators then
turned speech and notes into circulating books. A work may be authentically
Origenian without being the untouched product of one pen.

Survival added more collaborators. Greek copyists selected some works and let
others disappear. Rufinus translated homilies and \emph{On First Principles}
into Latin while openly smoothing passages he believed opponents had corrupted
or that conflicted with Origen elsewhere. Jerome translated other homilies,
quoted Greek passages against Rufinus, and changed from promoter to prosecutor
of Origen's legacy. Catena compilers dismembered commentaries into attributed
excerpts. Every theological claim must therefore identify not only work and
locus but Greek text, Latin translation, fragment, excerpt, or hostile report.

\subsection{The comparative Bible project}

The name \emph{Hexapla} refers to Origen's monumental comparison of biblical
texts. Eusebius describes Hebrew, the Hebrew transliterated into Greek letters,
and Greek versions associated with Aquila, Symmachus, the Septuagint, and
Theodotion, with additional anonymous Greek versions for some books
(\emph{Ecclesiastical History} 6.16--17). Origen also used critical signs
associated with Alexandrian philology to mark relationships among readings.

``Six columns'' is a useful name, not an adequate codicological description of
every biblical book. The order of stages, relation between Hexapla and
Tetrapla, treatment of extra versions, physical format, purpose, and degree of
revision remain disputed. Was the central goal a corrected Greek ecclesial
text, a synopsis for argument with Jewish readers, an archive of textual
difference, or several purposes developed across decades? The fragments do not
permit one answer to exclude all others.

The project was neither a modern critical edition nor a crude search for one
original wording. It put difference where a reader could see it. Its great
size also made it vulnerable: the complete work was not normally copied, and
later readers often received extracts or a Septuagintal text marked from it.
Modern reconstruction, including the Hexapla Institute's critical project,
gathers dispersed Greek, Syriac, Latin, Armenian, and other witnesses. It
recovers readings, not Origen's desk in full.

Origen also reasoned explicitly about the scope and authorship of Scripture.
Eusebius 6.25 excerpts his list of the Hebrew books, his statement that the
Church receives four Gospels, and his judgment that Hebrews transmits Pauline
thought in a style and diction not simply Paul's own. The famous conclusion
that God alone knows who wrote Hebrews is thus an Origenian judgment preserved
by Eusebius, not a modern slogan detached from his canon work. These passages
show textual comparison serving ecclesial reception; they do not amount to one
unchanging canon list for every congregation and every category of book.

\subsection{What survives, and in what form}

\begingroup
\small
\setlength{\tabcolsep}{4pt}
\begin{longtable}{>{\raggedright\arraybackslash}p{0.20\linewidth}
                  >{\raggedright\arraybackslash}p{0.17\linewidth}
                  >{\raggedright\arraybackslash}p{0.27\linewidth}
                  >{\raggedright\arraybackslash}p{0.28\linewidth}}
\toprule
\textbf{Work or family} & \textbf{Setting} & \textbf{Principal survival} &
\textbf{Historical control}\\
\midrule
\endhead
\emph{On First Principles} & Alexandria, probably 220s & Four books mainly in
Rufinus's Latin; substantial Greek excerpts and hostile citations & An unusually
large early Christian systematic investigation. Rufinus's announced
correction and the polemical Greek fragments must be compared; neither alone is
a transparent autograph.\\
\emph{Commentary on John} & Begun at Alexandria, continued at Caesarea &
Substantial Greek books and fragments, fewer than the ancient total & A
long-running anti-Gnostic and exegetical work whose books need not express one
unchanged moment.\\
Genesis and Matthew commentaries & Both centers & Uneven Greek books and
sections, catena fragments, and Latin translations & Large commentaries differ
from homilies and scholia; a detached fragment requires attribution and context
before use.\\
\emph{Commentary on the Song of Songs} & Caesarean period, building on earlier
homilies & Principally Rufinus's Latin translation, with Greek fragments & Its
ecclesial and nuptial exegesis became especially influential in Latin spiritual
reading; the translation layer must remain visible.\\
\emph{Commentary on Romans} & Caesarean period & Rufinus's Latin version in ten
books, abridging and reshaping a larger Greek work; Greek fragments & A major
witness to Origen on Paul, freedom, law, Israel, and justification, but not a
transparent copy of the lost Greek commentary.\\
Homilies & Principally presbyteral preaching at Caesarea & Mostly Latin
translations by Rufinus and Jerome; Greek fragments; twenty-nine Greek Psalm
homilies identified in 2012 & Congregational speech passed through stenography,
editing, translation, and collection. The Munich find materially enlarges the
controlled Greek homiletic voice.\\
\emph{On Prayer} and \emph{Exhortation to Martyrdom} & Caesarean spiritual
instruction, different occasions & Greek texts survive & Direct witnesses to
prayer, freedom, temptation, and confession; neither is a complete summary of
Origen's theology.\\
\emph{On Resurrection} & Early career; exact stages disputed & Lost as a whole;
fragments and reports, including hostile reception & The title and dispute are
secure, but later accusations cannot be treated as a continuous surviving
treatise.\\
\emph{On Pascha} & Date disputed & Greek papyrus portions from Tura & A material
recovery of paschal exegesis whose damaged state limits reconstruction.\\
\emph{Against Celsus} & Late, about 248 & Eight Greek books survive & Origen
preserves much of the lost opponent while arranging the answer; quotation does
not give unmediated access to Celsus.\\
\emph{Dialogue with Heraclides} & Ecclesial consultation, probably 240s & Greek
papyrus discovered at Tura in 1941 & A transcript-like record of oral exchange,
not necessarily an author-revised treatise or a council decree.\\
Letters & Across the career & A few complete letters or exchanges and many
extracts or notices & Eusebius saw a much larger correspondence. Surviving
pieces cannot make a complete network or development diary.\\
Hexaplar materials & Begun by the Alexandrian period, developed at Caesarea &
Dispersed fragments and daughter traditions & Layout, stages, and intention
remain active critical questions; ``the six-column Bible'' is shorthand.\\
\emph{Philocalia of Origen} & Fourth-century reception & Greek anthology of
extracts, traditionally associated with Basil and Gregory of Nazianzus & A
major route of authentic Origenian material and later selection, not a book
compiled by Origen himself.\\
\bottomrule
\end{longtable}
\endgroup

\subsection{Two discoveries that changed the modern author}

In 1941 workers at Tura, south of Cairo, uncovered a cache of papyrus codices.
The Biblioth\`eque nationale de France authority record describes approximately
two thousand pages, chiefly works of Origen and Didymus the Blind. Among the
Origen material were the
\emph{Dialogue with Heraclides} and Greek portions of works otherwise poorly
preserved. The find restored a setting in which bishops question one another
and formulations are tested aloud. It complicated the image built only from
systematic treatises and posthumous accusation.

In 2012 Marina Molin Pradel recognized that a twelfth-century Munich
manuscript, Bayerische Staatsbibliothek Cod. graec. 314, contained twenty-nine
homilies on Psalms corresponding to works in
Jerome's ancient list. Lorenzo Perrone and the editorial team tested the
attribution against language, parallel Latin and catena material, and the
historical corpus before publication. ``Rediscovery'' here means scholarly
identification of a manuscript already held in a library, not recovery of an
autograph. The event is nevertheless large: substantial Greek sermons now
stand beside the Latin translations through which Origen the preacher had
principally been known.

These discoveries expose the contingency of reputation. The Origen condemned
through propositions, the Origen corrected by translators, and the Origen
abstracted into spiritual exegesis were all partial. Papyrus and manuscript
work restored a consultant and preacher without making the contested thinker
disappear.
